\clearpage\pagestyle{otMain}\setlength{\otBlocEspacement}{2.5mm}\fontsize{11}{15}\selectfont\corrigesChapitre{16}{Triangles et quadrilatères}
\begin{blocLivre}{cor-16-1}\begin{corrige}{1}{Reconnaître les triangles}\demarche{Lire les codes avant de juger à l’œil.}
Le triangle 1 porte un petit carré : $\res{\text{rectangle}}$. Le 2 a deux côtés codés égaux : $\res{\text{isocèle}}$. Le 3 a trois côtés codés égaux : $\res{\text{équilatéral}}$. Le 4 n’a aucun code : $\res{\text{quelconque}}$.
\label{fig-cor-16-1}\figureGeo[corrige]{tri-1}
\end{corrige}\end{blocLivre}
\begin{blocLivre}{cor-16-2}\begin{corrige}{2}{Coder une figure}\demarche{Coder uniquement ce qui est donné.}
On code [MN] et [MP] avec le même petit trait : le triangle est $\res{\text{isocèle}}$ de sommet principal M. Le corrigé montre les codes en vert.
\label{fig-cor-16-2}\figureGeo[corrige]{tri-2}
\end{corrige}\end{blocLivre}
\begin{blocLivre}{cor-16-3}\begin{corrige}{3}{Le triangle équilatéral}\demarche{Garder la même ouverture de compas pour les deux arcs.}
On trace [BC] de 4 cm, puis deux arcs de rayon 4 cm centrés en B et en C ; leur croisement donne A. Les trois côtés mesurent 4 cm : les arcs du corrigé, en vert, montrent la méthode.
\label{fig-cor-16-3}\figureGeo[corrige]{tri-3}
\end{corrige}\end{blocLivre}
\begin{blocLivre}{cor-16-4}\begin{corrige}{4}{Reconnaître les quadrilatères}\demarche{Vérifier angles puis côtés.}
Quatre angles droits et côtés opposés égaux : $\res{\text{rectangle}}$. Quatre côtés égaux sans angle droit : $\res{\text{losange}}$. Quatre angles droits et quatre côtés égaux : $\res{\text{carré}}$. Le dernier, sans codes : $\res{\text{quadrilatère quelconque}}$.
\label{fig-cor-16-4}\figureGeo[corrige]{tri-4}
\end{corrige}\end{blocLivre}
\begin{blocLivre}{cor-16-5}\begin{corrige}{5}{Carré ou losange ?}\demarche{Tester un angle à l’équerre avant de conclure.}
$\res{\text{Non}}$ : quatre côtés égaux font un losange. Sans angle droit vérifié à l’équerre, rien ne prouve que c’est un carré. Ici, l’équerre ne s’ajuste pas : c’est un losange non carré.
\label{fig-cor-16-5}\figureGeo[corrige]{tri-5}
\end{corrige}\end{blocLivre}
\begin{blocLivre}{cor-16-6}\begin{corrige}{6}{Le triangle rectangle}\demarche{Commencer par l’angle droit, finir par l’hypoténuse.}
On trace l’angle droit en A à l’équerre, puis AB $=5$ cm et AC $=3$ cm sur ses côtés, et on joint B à C. Sur une construction soignée, on lit BC $\approx\res{5{,}8\ \mathrm{cm}}$.
\label{fig-cor-16-6}\figureGeo[corrige]{tri-6}
\end{corrige}\end{blocLivre}
\begin{blocLivre}{cor-16-7}\begin{corrige}{7}{Le triangle isocèle}\demarche{Les deux arcs portent la même ouverture de compas.}
On trace [EF] de 4 cm, puis deux arcs de rayon 6 cm centrés en E et en F : leur croisement donne D. Les codes affirment DE $=$ DF $=6$ cm.
\label{fig-cor-16-7}\figureGeo[corrige]{tri-7}
\end{corrige}\end{blocLivre}
\begin{blocLivre}{cor-16-8}\begin{corrige}{8}{Le rectangle aux instruments}\demarche{Élever les perpendiculaires aux deux bouts du premier côté.}
Perpendiculaires en A et en B, report de 4 cm, puis jonction. L’équerre confirme le quatrième angle $\res{\text{droit}}$ : la construction est correcte si les trois angles tracés le sont.
\label{fig-cor-16-8}\figureGeo[corrige]{tri-8}
\end{corrige}\end{blocLivre}
\begin{blocLivre}{cor-16-9}\begin{corrige}{9}{Le carré penché}\demarche{S’appuyer sur le côté donné, pas sur les bords de la feuille.}
La méthode ne change pas : perpendiculaires aux extrémités du côté donné, report de $4{,}5$ cm, jonction. Le carré reste un carré, $\res{\text{quelle que soit son orientation}}$.
\label{fig-cor-16-9}\figureGeo[corrige]{tri-9}
\end{corrige}\end{blocLivre}
\begin{blocLivre}{cor-16-10}\begin{corrige}{10}{Le losange au compas}\demarche{Le compas garde la même ouverture du début à la fin.}
Depuis A, reporter 5 cm sur les deux côtés de l’angle : B et D. Puis deux arcs de rayon 5 cm centrés en B et en D donnent C. Les quatre côtés, codés, mesurent $\res5$ cm.
\label{fig-cor-16-10}\figureGeo[corrige]{tri-10}
\end{corrige}\end{blocLivre}
\begin{blocLivre}{cor-16-11}\begin{corrige}{11}{Une figure composée}\demarche{Placer tous les points nommés avant les derniers tracés.}
Le rectangle se construit comme à l’exercice 8 ; M est à 3 cm de E. Les segments [MG] et [MH] partagent le rectangle en trois triangles. MG $=$ MH par symétrie de la figure : le triangle MGH est $\res{\text{isocèle}}$.
\label{fig-cor-16-11}\figureGeo[corrige]{tri-11}
\end{corrige}\end{blocLivre}
\begin{blocLivre}{cor-16-12}\begin{corrige}{12}{Le programme en désordre}\demarche{Chercher ce dont chaque étape a besoin.}
L’ordre correct est $\res{\text{b, d, a, c}}$ : le segment d’abord, le réglage du compas, les arcs qui créent A, puis les côtés. Un programme se lit et s’exécute dans l’ordre.
\label{fig-cor-16-12}\figureGeo[corrige]{tri-12}
\end{corrige}\end{blocLivre}
\begin{blocLivre}{cor-16-13}\begin{corrige}{13}{La figure qui ment}\demarche{Confronter l’affirmation aux instruments.}
Un seul angle non droit suffit : ce n’est $\res{\text{pas un rectangle}}$. La figure garde quatre côtés : c’est un $\res{\text{quadrilatère}}$. Une annonce ou une apparence ne remplace pas la vérification.
\label{fig-cor-16-13}\figureGeo[corrige]{tri-13}
\end{corrige}\end{blocLivre}
\begin{blocLivre}{cor-16-14}\begin{corrige}{14}{Le plus grand tour}\demarche{Compter les côtés avant de multiplier.}
Carré : $4\times6=\res{24}$ cm. Triangle : $3\times6=\res{18}$ cm. Le $\res{\text{carré}}$ a le plus grand tour : il a un côté de plus, de même longueur.
\label{fig-cor-16-14}\figureGeo[corrige]{tri-14}
\end{corrige}\end{blocLivre}
\begin{blocLivre}{cor-16-15}\begin{corrige}{15}{Les diagonales du rectangle}\demarche{Mesurer les deux diagonales avec soin.}
Les deux diagonales mesurent chacune $\res{10\ \mathrm{cm}}$ : dans un rectangle, les diagonales ont la $\res{\text{même longueur}}$ et se coupent en leur milieu. Cette observation se vérifie à la règle sur toute construction soignée.
\label{fig-cor-16-15}\figureGeo[corrige]{tri-15}
\end{corrige}\end{blocLivre}
